\section{Absorbing four parity conditions into the physical anchor}
\label{sec:packed-anchor-note}

The construction below removes the sixteen private rational coordinates
used by the four odd-index charts in the Version~10 source.  It adds no
coordinate: their arithmetic content is absorbed into the existing
three-witness physical anchor.  This is an equivalence of the rational
source systems on their unchanged point domains.  It does not constitute
an independent verification of the source's global Diophantine-definition
theorem.

\paragraph{Provenance and inherited inputs.}
The baseline is
\href{https://github.com/elliotglazer/diophantine-z-q/blob/036cebe4f1a4aa7ca89820a9b48f3629b4c364d6/manuscripts/quartic_v10.tex}{\emph{Z is Diophantine in Q}, Version~10},
at commit
\texttt{036cebe4f1a4aa7ca89820a9b48f3629b4c364d6}.
We retain its point-domain and group statements, its integral formal-group
construction, and its three-witness quaternion criterion.  Their precise
source labels are listed below.  The new ingredients proved here are the
leading-residue calculation for arbitrary multiples of the fixed point,
the resulting product parity test, and its combination with the anchor by
separating even and odd valuations.  No height estimate or prime-pattern
argument is used in proving this replacement.

\subsection{The fixed curve and the residue calculation}

Put
\[
 \rho=\sqrt2,\qquad K=\mathbb Q(\rho),\qquad
 D=11-6\rho,\qquad E:y^2=x^3-Dx,
\]
\[
 P=\left(\frac{57}{8}-\frac72\rho,
          -\frac{165}{8}+\frac{417}{32}\rho\right).
\]
The retained source statements give $2E(K)=\mathbb ZP$, with $P$ of
infinite order.  Each of the four root points is a finite nonidentity
point $Q_i=[n_i]P$, where $n_i\in\mathbb Z\setminus\{0\}$.  The existing
square-abscissa charts are retained.  Write
\[
 C_{\mathrm{prod}}=\prod_{i=1}^4x(Q_i)=C_0+C_1\rho,
 \qquad C_0,C_1\in\mathbb Q.
\]
This product and the scale $s_0=C_0^2$ are already materialized in the
source.

At the unique place above $3$, use $v(3)=1$; the residue field is
$\mathbb F_9=\mathbb F_3(\bar\rho)$, with $\bar\rho^2=2$.
Reduction of the displayed point gives
$\bar P=(\bar\rho,0)$, of order two, on the good reduction of $E$.
The formal parameter is $z=-x/y$.

\paragraph{Lemma 1 (leading residues).}
For every nonzero integer $n$, define
\[
 e(n)=\begin{cases}
 0,&n\text{ odd},\\
 2+2v_3(n),&n\text{ even}.
 \end{cases}
\]
Then $3^{e(n)}x([n]P)$ is integral at the place above $3$, and
\begin{equation}\label{eq:pa-leading-residues}
 3^{e(n)}x([n]P)\equiv
 \begin{cases}
 \rho,&n\text{ odd},\\
 2\rho,&n\text{ even}
 \end{cases}\pmod3.
\end{equation}

\begin{proof}
For odd $n$, the reduction of $[n]P$ is $\bar P$, so the first assertion
follows immediately.  For even $n$, write
$n=2\cdot3^ku$, where $k\ge0$ and $u\in\mathbb Z$ is not divisible by
$3$.  Direct duplication of the displayed $P$ gives
\[
 z(2P)=
 \frac{3396356159088-4135204317756\rho}{3287811625151},
 \qquad \frac{z(2P)}3\equiv2+\rho\pmod3.
\]
The denominator is a $3$-adic unit, and each numerator coefficient has
valuation one.

The integral formal-group multiplication series are odd, with
$[m](z)=mz+O(z^3)$ for every integer $m$.  At depth $r=v(z)\ge1$,
multiplication by a $3$-adic unit integer preserves depth and multiplies
the normalized leading residue by that integer.  Multiplication by $3$
has depth $r+1$: after division by $3^{r+1}$ its error has valuation at
least $3r-(r+1)=2r-1\ge1$.  Its normalized leading residue is therefore
unchanged.  Induction, with $r=k+1$, gives
\[
 3^{-r}z([n]P)\equiv u(2+\rho)\pmod3.
\]
On the formal branch, $xz^2\equiv1\pmod3$.  One may see this directly
from $xz^2=x^2/(x^2-D)$ and $v(x)<0$.  Since $u^2=1$ in
$\mathbb F_3^*$ and
$(2+\bar\rho)^{-2}=2\bar\rho$ in $\mathbb F_9$, the even case of
\eqref{eq:pa-leading-residues} follows.  Negative $u$ is allowed in the
multiplication series; equivalently, inversion changes $z$ to $-z$ and
leaves $x$ unchanged.  Thus negative $n$ is covered as well.
\end{proof}

\paragraph{Corollary 2 (one coefficient detects all four parities).}
Let $E_0=\sum_i e(n_i)$.  Then
\begin{equation}\label{eq:pa-product-parity}
 v_3(C_0)=-E_0,\qquad
 C_0\in\mathbb Z_{(3)}
 \quad\Longleftrightarrow\quad n_1,n_2,n_3,n_4\text{ are all odd}.
\end{equation}
In particular $C_0\ne0$.

\begin{proof}
If $j$ of the four labels are even, multiplying
\eqref{eq:pa-leading-residues} gives
\[
 3^{E_0}C_{\mathrm{prod}}\equiv
 2^j\rho^4=2^j\pmod3.
\]
Its residue is a nonzero element of $\mathbb F_3$.  In the integral
basis $(1,\rho)$, the rational coefficient $3^{E_0}C_0$ is consequently
a $3$-adic unit.  This proves the equality of valuations, and every
$e(n_i)$ is nonnegative and vanishes exactly for odd $n_i$.
\end{proof}

\subsection{The place above two and the packed equation}

At the place above $2$, normalize $v(2)=1$, so $v(\rho)=1/2$.
For rational $a,b$, the two possible finite valuations have different
fractional parts; hence
\begin{equation}\label{eq:pa-two-basis}
 v(a+b\rho)=\min\{v_2(a),v_2(b)+\tfrac12\}.
\end{equation}
The displayed point has $v(z(P))=3/2$ and belongs to the integral
formal subgroup at infinity, without requiring a minimal model.
All its integer multiples remain in that subgroup.  For odd $n$, the
unit linear coefficient of $[n](z)$ dominates the higher integral terms,
so $v(z([n]P))=3/2$ and $v(x([n]P))=-3$.
If all four labels are odd, it follows that
$v(C_{\mathrm{prod}})=-12$; by \eqref{eq:pa-two-basis}, in particular
$v_2(C_0)\ge-12$.  Therefore
\begin{equation}\label{eq:pa-B}
 B=2^{12}C_0=4096C_0\in\mathbb Z_{(2)}
 \quad\text{whenever all four labels are odd}.
\end{equation}

Let $Q_2$ be the root called $S$ in the source, and let $R=Q_2-P$.
The unchanged subtraction guard makes $R$ a finite nonidentity point,
so $R$ is a nonzero integer multiple of $P$.  Write
\[
 f(R)=\frac1{x(R)}=f_0+f_1\rho,\qquad
 C_{\mathrm a}=3^{2^{25}},\qquad
 \mathsf S=C_{\mathrm a}^2=3^{2^{26}}.
\]
Formal-group membership of $R$ gives $v(f(R))>0$ above $2$, and
\eqref{eq:pa-two-basis} shows that $f_0,f_1\in\mathbb Z_{(2)}$.
Since $C_{\mathrm a}$ is odd,
$f_0/C_{\mathrm a}$ and $f_1/C_{\mathrm a}$ are always $2$-integral,
without assuming the anchor or any root parity condition.

Set $q_2(a,b)=a^2+ab-b^2$.  The old physical anchor materializes
$A_{\mathrm{phys}}=q_2(f_0,f_1)$ and asks for rational $u,v,w$ with
\[
 u^2+6v^2-2w^2-3A_{\mathrm{phys}}^2+2\mathsf S^2=0.
\]
Replace its materialized coordinate by
\begin{equation}\label{eq:pa-Theta}
 \Theta=q_2(f_0,f_1)+3\mathsf S B^2
        =q_2(f_0,f_1)+50331648\,\mathsf S s_0,
 \qquad s_0=C_0^2,
\end{equation}
and retain the same three-witness equation:
\begin{equation}\label{eq:pa-anchor}
 u^2+6v^2-2w^2-3\Theta^2+2\mathsf S^2=0.
\end{equation}

\paragraph{Proposition 3 (exact anchor equivalence).}
On the unchanged point domains, equations
\eqref{eq:pa-Theta}--\eqref{eq:pa-anchor} have rational witnesses if and
only if the old physical anchor has rational witnesses and all four
root labels are odd.

\begin{proof}
Use the source's quaternion criterion, valid for every $t\in\mathbb Q$:
\begin{equation}\label{eq:pa-quaternion}
 t\in\mathbb Z_{(2)}\cap\mathbb Z_{(3)}
 \ \Longleftrightarrow\
 \exists u,v,w\in\mathbb Q:\quad
 u^2+6v^2-2w^2-3t^2+2=0.
\end{equation}
It comes from the quadratic embedding criterion for $(2,3)_{\mathbb Q}$,
ramified exactly at $2$ and $3$; the source includes $t=0$ explicitly.
Dividing \eqref{eq:pa-anchor} by $\mathsf S^2$ applies this criterion to
\[
 t=\Theta/\mathsf S=q_2(a,b)+3B^2,
 \qquad a=f_0/C_{\mathrm a},\quad b=f_1/C_{\mathrm a}.
\]
The form $q_2$ is anisotropic modulo both $2$ and $3$.
In particular
$v_3(q_2(a,b))=2\min\{v_3(a),v_3(b)\}$.
This valuation is even, whereas $v_3(3B^2)=1+2v_3(B)$ is odd;
they cannot cancel.  Thus
\[
 t\in\mathbb Z_{(3)}
 \quad\Longleftrightarrow\quad
 a,b,B\in\mathbb Z_{(3)}.
\]
The possibility $q_2(a,b)=0$ causes no exception (use $v(0)=+\infty$).
By \eqref{eq:pa-product-parity}, the condition on $B$ is exactly
oddness of all four labels.  When it holds, \eqref{eq:pa-B} and the
preceding automatic integrality of $a,b$ imply $t\in\mathbb Z_{(2)}$.

Finally, applying \eqref{eq:pa-quaternion} to the old anchor, and using
$v_p(q_2(a,b))=2\min(v_p(a),v_p(b))$ for $p=2,3$, identifies that anchor
with $a,b\in\mathbb Z_{(2)}\cap\mathbb Z_{(3)}$.
Their $2$-integrality is automatic here.  Hence the two descriptions
are equivalent in both directions, including existence of the old
or new conic witnesses as required.
\end{proof}

\subsection{Projection, quadratic degree, and source dependencies}

Delete the four old odd-chart augmentations, each with four private
rational coordinates.  The surviving charts still describe the same
points $Q_i=(h_i^2,h_ij_i)$.  For each odd multiple the source's odd-chart
proposition supplies an augmented representation.  If its $h_i,j_i$
differ from the current representation, $h_i^2=x(Q_i)\ne0$ implies
that they are obtained by the simultaneous switch
$(h_i,j_i)\mapsto(-h_i,-j_i)$.  This preserves $x$, $y=hj$, and
$z=-h/j$.  In the pinned source, the remaining point functions and their
cleared-denominator equations are invariant under this switch; the
subtraction slopes use $hj$, and the chart parameter uses $-h/j$.
Thus old odd-chart witnesses can be restored after these independent
sign switches.  Proposition~3 restores the old anchor witnesses as
well.  Conversely, any old solution satisfies the parity conditions
and extends to the new anchor.  The existential projection to the
original input is unchanged; no unique correspondence between witness
tuples is asserted.

Equation~\eqref{eq:pa-Theta} must use the existing coordinate
$s_0=C_0^2$, and $\Theta$ must remain materialized.  It is quadratic
even when $\mathsf S$ is represented by the last coordinate of the
existing constant chain.  The conic \eqref{eq:pa-anchor} is also
quadratic.  Consequently the usual sum of squared residuals still has
degree at most four.  The arithmetic replacement removes precisely
sixteen private coordinates and adds none; any further compiler
eliminations and final coefficient statistics are separate finite
calculations.

The inherited results used above are identified by these literal
labels in the pinned manuscript:
\begin{itemize}
 \item \texttt{subsec:v5-point-source}, \texttt{lem:v5-point-domains},
 and \texttt{v9chart:prop:chart}: the normalized group, point domains,
 and square-abscissa charts;
 \item \texttt{subsec:o8-anchor-subgroup}: the integral formal group,
 the exact value of $z(2P)$, and the local calculations at $2$ and $3$;
 \item \texttt{lem:v5-three-witness-anchor} and
 \texttt{lem:v9-bridge}: the quaternion criterion and Version~9
 anchor constants;
 \item \texttt{v9:prop:odd-chart}: existential recovery of the removed
 odd-chart witnesses.
\end{itemize}
The projection argument uses none of the claimed global height or
prime-selection theorems.  Those remain inherited dependencies when
this source-system equivalence is used to produce a smaller claimed
defining polynomial for $\mathbb Z$ in $\mathbb Q$.
